\section{Context Engineering: hacer que la IA entienda tu código}

\subsection{De un prompt simple a un documento de especificación al estilo de un gerente de producto}

Para cambios simples, basta con un prompt improvisado. Pero para tareas complejas, el desarrollador necesita desempeñar el papel de un gerente de producto (PM) y redactar un documento de especificación (spec doc) cuidadosamente diseñado.

\begin{importantbox}{Estructura completa de un spec doc}
\begin{itemize}
    \item \textbf{Goal}: cuál es el propósito del cambio
    \item \textbf{Definitions}: el conocimiento previo que el LLM necesita conocer
    \item \textbf{Plan}: el desglose de alto nivel del plan de implementación
    \item \textbf{Source files}: los archivos de código involucrados y su relevancia
    \item \textbf{Test cases}: cómo realizar las pruebas
    \item \textbf{Edge cases}: los casos especiales que hay que manejar
    \item \textbf{Out-of-scope}: qué contenido no debe modificarse
    \item \textbf{Extensions}: las posibles direcciones de extensión futuras (ayuda al LLM a hacer un diseño con visión de futuro y evitar atajos)
\end{itemize}
\end{importantbox}

\subsection{Optimizar la base de código para la IA}

\begin{warningbox}{El origen de la confusión del LLM suele ser la base de código misma}
Muchas veces el LLM produce malos resultados no porque la capacidad del modelo sea insuficiente, sino porque la base de código está demasiado desordenada. Si un ingeniero humano recién incorporado al equipo también se confundiría con tu base de código, el LLM se confundirá igual.
\end{warningbox}

Las dimensiones clave para optimizar la base de código:
\begin{itemize}
    \item \textbf{Repo orientation}: una introducción general del repositorio
    \item \textbf{File structure}: una estructura de directorios clara
    \item \textbf{Setup and environment}: las instrucciones de configuración del entorno
    \item \textbf{Code style}: las convenciones de estilo de código
    \item \textbf{Access patterns}: los patrones de acceso a datos
    \item \textbf{APIs and contracts}: las interfaces de API y sus convenciones
\end{itemize}

Nota: \textbf{el diseño de repositorio único} se recomienda ampliamente en los escenarios de programación con IA, porque permite que el Agent entienda con mayor facilidad el panorama completo del proyecto.

\subsection{Archivos de configuración del Agent}

\begin{itemize}
    \item \textbf{CLAUDE.md}: el archivo de descripción del proyecto que Claude Code carga automáticamente; es adecuado para colocar los comandos de bash de uso frecuente, los archivos centrales y las funciones de utilidad, la guía de estilo de código, las instrucciones de prueba, etc.
    \item \textbf{.cursorrules}: el archivo de reglas del proyecto de Cursor
    \item \textbf{AGENTS.md}: un archivo de formato abierto para las instrucciones del Agent
    \item \textbf{llms.txt}: proporciona una guía de navegación para los rastreadores web de LLM
\end{itemize}

\begin{warningbox}{El Agent no siempre respeta los archivos de configuración}
Estas descripciones e instrucciones son solo una \textbf{guía} (guidance); el Agent no siempre las respeta de manera estricta. No dependas de los archivos de configuración para implementar restricciones de seguridad o invariantes críticos.
\end{warningbox}

\subsection{Resumen de la sección}

El núcleo del context engineering es: proporcionar al LLM la información de contexto óptima. Esto incluye redactar un spec doc estructurado, optimizar la organización de la base de código y aprovechar los archivos de configuración del Agent. El objetivo es que tanto los humanos como el Agent puedan entender el proyecto con facilidad.
